% !TeX root = slides.tex
\documentclass[aspectratio=169,10pt]{beamer}
\PassOptionsToPackage{ngerman}{babel}
\usepackage[T1]{fontenc}
\usepackage[utf8]{inputenc}
\usepackage[english]{babel}
\usepackage{lmodern}
\usepackage{booktabs}
\usepackage{array}
\usepackage{tikz}
\usetikzlibrary{arrows.meta,positioning,calc}

\definecolor{Ink}{HTML}{172B43}
\definecolor{Muted}{HTML}{526579}
\definecolor{Teal}{HTML}{087F8C}
\definecolor{Indigo}{HTML}{4956A3}
\definecolor{Amber}{HTML}{DD8D29}
\definecolor{Rust}{HTML}{AF4D35}
\definecolor{Pale}{HTML}{EFF4F7}
\definecolor{Line}{HTML}{DBE4EB}
\setbeamercolor{normal text}{fg=Ink,bg=white}
\setbeamercolor{structure}{fg=Teal}
\setbeamercolor{frametitle}{fg=Ink}
\setbeamercolor{alerted text}{fg=Rust}
\setbeamercolor{block title}{fg=Ink,bg=Pale}
\setbeamercolor{block body}{fg=Ink,bg=Pale}
\setbeamerfont{frametitle}{size=\Large,series=\bfseries}
\setbeamerfont{block title}{size=\normalsize,series=\bfseries}
\setbeamersize{text margin left=8mm,text margin right=8mm}
\setbeamertemplate{navigation symbols}{}
\setbeamertemplate{itemize item}{\color{Teal}\small\textbullet}
\setbeamertemplate{frametitle}{\vspace{3mm}\insertframetitle\par\vspace{1mm}\textcolor{Line}{\rule{\textwidth}{.5pt}}}
\setbeamertemplate{footline}{%
  \hspace{8mm}\textcolor{Muted}{\scriptsize SSBI · Task 3 · 20,000 PBMCs}\hfill
  \textcolor{Muted}{\scriptsize\insertframenumber}\hspace{8mm}\vspace{3mm}}
\setlength{\parskip}{2pt}
\newcommand{\takeaway}[1]{\vspace{2mm}\begin{beamercolorbox}[sep=1.7mm,wd=\linewidth]{block body}\small\textbf{#1}\end{beamercolorbox}}
\newcommand{\source}[1]{\vfill{\fontsize{6.5}{8}\selectfont\color{Muted}#1\par}}
\newcommand{\up}{\ensuremath{\uparrow}}
\newcommand{\down}{\ensuremath{\downarrow}}

\usepackage{amsmath,amssymb}
\AtBeginDocument{\selectlanguage{ngerman}}
\graphicspath{{figures/}}
\newcommand{\PCOne}{12.60}
\newcommand{\PCTwo}{7.69}
\newcommand{\PCTogether}{20.29}
\newcommand{\PCSelected}{91.50}

\hypersetup{pdftitle={CellCNN-Varianten: 100 gemeinsame Spendersplits},pdfauthor={SSBI-Gruppenprojekt},colorlinks=true,urlcolor=Teal}
\setbeamertemplate{footline}{%
  \hspace{8mm}\textcolor{Muted}{\scriptsize SSBI · Aufgabe 6 · 100 Splits · 20 Spender}\hfill
  \textcolor{Muted}{\scriptsize\insertframenumber}\hspace{8mm}\vspace{3mm}}
\newcommand{\eyebrow}[1]{{\color{Teal}\bfseries #1}\par\smallskip}
\newcommand{\plotcontext}[1]{\par{\scriptsize\color{Muted}#1}\par}

\begin{document}

\begin{frame}{Vier Zellfilter-Modelle im gemeinsamen Vergleich}
  \eyebrow{Gleiche Daten, gleiche Spendersplits, gleiche Half-Max-Regel}
  \begin{tabular}{@{}p{.25\textwidth}p{.47\textwidth}p{.20\textwidth}@{}}\toprule
    Modell & Zellresponse $h_k(x)$ & Pooling\\\midrule
    \textcolor{Teal}{CellCNN} & $\operatorname{ReLU}(b_k+w_k^\top x)$ & hartes Top-1\,\%\\[2mm]
    \textcolor{Indigo}{Quadratic Top-1\,\%} & $\operatorname{ReLU}(b_k+w_k^\top x+\sum_j q_{kj}x_j^2)$ & hartes Top-1\,\%\\[2mm]
    \textcolor{Amber}{Quadratic Soft-$\alpha$} & $\operatorname{ReLU}(b_k+w_k^\top x+\sum_j q_{kj}x_j^2)$ & lernbares Soft-$\alpha$\\[2mm]
    \textcolor{Rust}{Prototyp Soft-$\alpha$} & $\operatorname{ReLU}(\rho_k-\frac1d\sum_j a_{kj}(x_j-c_{kj})^2)$ & lernbares Soft-$\alpha$\\\bottomrule
  \end{tabular}
  \medskip
  Alle Modelle aggregieren Zellresponses und erzeugen danach lineare Klassenlogits.
  \takeaway{Netzwerk-Klassifikation und Assoziation der ausgewählten Population werden getrennt bewertet.}
\end{frame}

\begin{frame}{100 Auswertungen, aber nur 20 unabhängige Spender}
  \begin{columns}[T,onlytextwidth]
    \begin{column}{.47\textwidth}
      \eyebrow{Gemeinsame Datengrundlage}
      \begin{tabular}{@{}lr@{}}\toprule
      Unabhängige Spender & 20\\
      CMV$-$ / CMV$+$ & 11 / 9\\
      Marker & 37\\
      Gate & \texttt{gated\_alive}\\
      Äußere Train-/Testspender & 14 / 6\\\bottomrule
      \end{tabular}
      \medskip\par
      \small Transformation $\operatorname{asinh}(x/5)$; Scaler ausschließlich auf Inner-Train-Spendern.
    \end{column}
    \begin{column}{.49\textwidth}
      \eyebrow{Dasselbe Auswahlverfahren}
      \begin{itemize}
        \item Drei Inner-Folds $\times$ drei Filterzahlen: neun Kandidaten je Split.
        \item Auswahl nach Validierungsgenauigkeit, AUC und Verlust.
        \item Ausgewähltes Inner-Modell wird ohne Refit getestet.
        \item 3.000 Zellen/Trainingsbag, 200 Bags/Spender; fünf Prediction-Bags mit bis zu 20.000 Zellen.
      \end{itemize}
    \end{column}
  \end{columns}
  \takeaway{600 Testauftritte je Methode sind keine 600 unabhängigen Spender. Die Analyse ist deskriptiv.}
\end{frame}

\begin{frame}{Eine gemeinsame Half-Max-Definition für alle vier Modelle}
  \begin{columns}[T,onlytextwidth]
    \begin{column}{.54\textwidth}
      \eyebrow{Filterauswahl und feste Trainingsschwelle}
      \[\Delta_k=V_{1k}-V_{0k},\qquad k^*=\arg\max_{k:\Delta_k>0}\Delta_k\]
      \[M=\max_{x\in\text{Inner-Train-Zellen}}h_{k^*}(x),\qquad t=\tfrac12M\]
      \[f_i=\frac{\#\{x\in\text{Teststichprobe }i:h_{k^*}(x)>t\}}{N_i}\]
      \small Gleiche feste Stichprobe von bis zu 20.000 Zellen je Testspender. Die strikte Regel gilt auch bei $M=0$.
    \end{column}
    \begin{column}{.42\textwidth}
      \eyebrow{Zwei Populationsmetriken}
      \[\mathrm{AUC}_{\mathrm{freq}}=\operatorname{AUC}(y_i,f_i)\]
      \[\mathrm{Effekt}=\bar f_{\mathrm{CMV+}}-\bar f_{\mathrm{CMV-}}\]
      \small Ohne positiven Filter bleiben die Werte fehlend. Kein nachträgliches Umdrehen der AUC.
      \medskip\par
      Für Frequency-Vergleiche: \textbf{\CommonFrequencySplits{} gemeinsam bestimmbare Splits}.
    \end{column}
  \end{columns}
  \takeaway{Half-Max ist eine Zellselektion. Es wird kein Clustering ausgeführt; CMV ist kein Zelltyp-Label.}
\end{frame}

\begin{frame}{Netzwerk-ROC auf denselben 100 Splits}
  \begin{columns}[c,onlytextwidth]
    \begin{column}{.65\textwidth}
      \includegraphics[width=\linewidth,height=49mm,keepaspectratio]{network_roc.pdf}
    \end{column}
    \begin{column}{.32\textwidth}
      \eyebrow{Mittelwert / Median}
      {\scriptsize\begin{tabular}{@{}lrr@{}}\toprule
Modell & Mittel & Median\\\midrule
CellCNN & 0,8075 & 0,8750\\
Quad.~Top-1\,\% & 0,9125 & 1,0000\\
Quad.~Soft-$\alpha$ & 0,9125 & 1,0000\\
Prototyp & 0,7375 & 0,7500\\
\bottomrule\end{tabular}
}
      \medskip\par\small ROC-Kurven je Split berechnen, auf gemeinsamer FPR-Achse interpolieren und mitteln.
    \end{column}
  \end{columns}
  \takeaway{\NetworkTakeaway}
  \plotcontext{Legenden-AUC: Mittel der ursprünglichen Split-AUCs. Kein Pooling aller 600 Testauftritte zu einer ROC.}
\end{frame}

\begin{frame}{Wie stark schwankt die Klassifikation zwischen Splits?}
  \relax
  {\centering
    \includegraphics[width=\textwidth,height=53mm,keepaspectratio]{network_metrics.pdf}
  \par}
  \takeaway{\ClassificationTakeaway}
  \plotcontext{ROC-AUC, Average Precision und Balanced Accuracy je Split; BA bei Score-Schwelle 0,5. Boxen Q1--Q3, Linie Median; Streuung ist kein Konfidenzintervall.}
\end{frame}

\begin{frame}{Frequency-ROC der ausgewählten Zellpopulation}
  \begin{columns}[c,onlytextwidth]
    \begin{column}{.65\textwidth}
      \includegraphics[width=\linewidth,height=49mm,keepaspectratio]{frequency_roc.pdf}
    \end{column}
    \begin{column}{.32\textwidth}
      \eyebrow{Gemeinsame Splits}
      {\scriptsize% Generated from the frozen benchmark; numerical values are unchanged.
\begin{tabular}{@{}lrr@{}}\toprule
Model & AUC & Effect*\\\midrule
CellCNN & 0.8559 & 0.740\\
Quad.\ Top-1\% & 0.8610 & 0.570\\
Quad.\ Soft-$\alpha$ & 0.8559 & 0.535\\
Prototype & 0.7245 & 0.300\\
\bottomrule\end{tabular}
\par\smallskip{\tiny *Percentage points; common splits.}
}
      \medskip\par\small Alle vier Methoden werden auf denselben \CommonFrequencySplits{} bestimmbaren Splits verglichen.
    \end{column}
  \end{columns}
  \takeaway{\FrequencyTakeaway}
  \plotcontext{Gemessen wird CMV-Assoziation der Häufigkeit; daraus folgt keine Zelltypvalidierung.}
\end{frame}

\begin{frame}{Frequency-AUC und Effektgröße beantworten verschiedene Fragen}
  \relax
  {\centering
    \includegraphics[width=.94\textwidth,height=46mm,keepaspectratio]{frequency_metrics.pdf}
  \par}
  \takeaway{\EffectTakeaway}
  \plotcontext{AUC: Rangordnung der Donoren. Effekt: Differenz ihrer mittleren Populationsfrequenzen; Einheit Prozentpunkte.}
\end{frame}

\begin{frame}{Gepaarte Differenzen zur linearen CellCNN-Baseline}
  \relax
  {\centering
    \includegraphics[width=.94\textwidth,height=49mm,keepaspectratio]{paired_differences.pdf}
  \par}
  \takeaway{\PairedTakeaway}
  \plotcontext{Differenzen werden innerhalb desselben äußeren Splits gebildet. Positive Werte bedeuten eine höhere AUC als CellCNN.}
\end{frame}

\begin{frame}{Welche Frequenzen haben die einzelnen Spender?}
  \relax
  {\centering
    \includegraphics[width=\textwidth,height=51mm,keepaspectratio]{donor_frequencies.pdf}
  \par}
  \takeaway{Jeder Spender zählt als ein Punkt; wiederholte Testauftritte werden zuerst je Spender gemittelt.}
  \plotcontext{Kreise: CMV$-$; Dreiecke: CMV$+$. Striche: Mittel der Spenderpunkte. Je Split kann ein anderer Filter die Population definieren.}
\end{frame}

\begin{frame}{Was lernt der zusätzliche Poolingparameter?}
  \begin{columns}[c,onlytextwidth]
    \begin{column}{.58\textwidth}
      \includegraphics[width=\linewidth,height=51mm,keepaspectratio]{learned_alphas.pdf}
    \end{column}
    \begin{column}{.39\textwidth}
      \eyebrow{Ranganteil statt Populationsfrequenz}
      \[\alpha_k=\sigma(\theta_k),\quad \alpha_k^{(0)}=0{,}01\]
      \[m_{rk}=\sigma\!\left(\frac{\alpha_k-q_r}{0{,}002}\right)\]
      \small Mit $q_r=(r-1/2)/N$. Die Temperatur bleibt fest. Ein weiches Alpha von 1\,\% ist nicht identisch mit hartem Top-1\,\%.
    \end{column}
  \end{columns}
  \takeaway{\AlphaTakeaway}
\end{frame}

\begin{frame}{Einordnung des kontrollierten Vergleichs}
  \begin{itemize}
    \item \textbf{Quadratic Hard vs. Soft:} gleiche Zellresponse und L2-Strafe; geändert wird die Poolingfunktion.
    \item \textbf{CellCNN vs. Quadratic:} zusätzlicher quadratischer Term und dessen L2-Strafe.
    \item \textbf{Prototyp vs. Quadratic:} andere Geometrie, Initialisierung und Regularisierung; kein isolierter Alpha-Vergleich.
    \item \textbf{Begrenzung:} Varianten wurden nach früheren Ergebnissen entwickelt. Mehr Splits erzeugen keine neuen unabhängigen Donoren.
  \end{itemize}
  \medskip
  \textbf{Regularisierung:} CellCNN auf $W,V$; Quadratic auf $W,Q,V$.
  Prototyp: Output-L2 plus $10^{-3}\operatorname{mean}((a-1)^2)$.
  \takeaway{\FinalTakeaway}
  \plotcontext{Alle 100 vorgesehenen Splits werden berichtet. Code, Checkpoints, Tabellen und Abbildungen sind gemeinsam dokumentiert.}
\end{frame}

\end{document}
